% GENERATED FILE. Edit canonical lessons, then run npm run generate:books.
% canonical-source-hash: fnv1a64:4aa552e59c8788fd
% canonical-lessons: ML-C418-siksa, ML-C418-siksikkuka, ML-C418-vilakku, ML-C418-nirodhikkuka, ML-C418-anuvadikkuka

\chapter{Punishment and Permission}
\label{ch:ml-punishment-and-permission}
\hlchaptermodality{418}

\begin{chapteropening}
\textbf{By the end of this chapter:} \emph{I can say punishment, to punish, a ban, to prohibit and to allow.}

Complete the last lesson of chapter 418: I can say punishment, to punish, a ban, to prohibit and to allow.
\end{chapteropening}

\section[\texorpdfstring{\ml{ശിക്ഷ} (śikṣa)}{ശിക്ഷ (śikṣa)}]{\ml{ശിക്ഷ} (śikṣa) — punishment; a sentence}

\label{lesson:ML-C418-siksa}



\begin{quote}
Before the new one: say the Malayalam for legal, then the Malayalam for illegal.
\end{quote}

\subsection*{You'll want to know: \ml{ശിക്ഷ}}
\textbf{\ml{ശിക്ഷ}} — \emph{śikṣa} — \textquotedblleft{}punishment; a sentence\textquotedblright{}.

\begin{quote}
\ml{തെറ്റിന്} \ml{അവന്} \ml{ചെറിയ} \ml{ഒരു} \ml{ശിക്ഷ} \ml{കിട്ടി}.
\end{quote}

\begin{quote}
\emph{teṟṟinŭ avanŭ ceṟiya oru śikṣa kiṭṭi.}
\end{quote}

\begin{quote}
He got a small punishment for the mistake.
\end{quote}

A punishment is \ml{കിട്ടുക}, received, or \ml{വിധിക്കുക}, passed by a court.

\begin{cousinweb}
\textbf{\ml{ശിക്ഷ}} (\emph{śikṣa}) is borrowed from Sanskrit \emph{śikṣā}, \textquotedblleft{}learning, study\textquotedblright{}, which later also came to mean \textquotedblleft{}punishment\textquotedblright{} (CDIAL 12432). The lesson and the penalty share one word.
\end{cousinweb}

\begin{grammarlens}[title={What you've built}]
One more word for giving reasons and opinions.
\end{grammarlens}

\subsection*{Your turn}
\begin{itemize}
\raggedright
  \item \emph{Say it:} \emph{śikṣa}
  \item \emph{Say it:} \emph{śikṣa}, once more
  \item \emph{Say it:} \emph{teṟṟinŭ avanŭ ceṟiya oru śikṣa kiṭṭi.} — He got a small punishment for the mistake.
  \item \emph{Recall:} say the Malayalam for legal, then the Malayalam for illegal, then say \emph{śikṣa} again
\end{itemize}

\begin{tcolorbox}[breakable,colback=teal!4,colframe=teal!35!black,title={Before you move on}]
What does \ml{ശിക്ഷ} mean? (\textquotedblleft{}Punishment; a sentence\textquotedblright{}.) Say it once more.
\end{tcolorbox}
\section[\texorpdfstring{\ml{ശിക്ഷിക്കുക} (śikṣikkuka)}{ശിക്ഷിക്കുക (śikṣikkuka)}]{\ml{ശിക്ഷിക്കുക} (śikṣikkuka) — to punish}

\label{lesson:ML-C418-siksikkuka}



\begin{quote}
Before the new one: say the Malayalam for illegal, then the Malayalam for punishment.
\end{quote}

\subsection*{You'll want to know: \ml{ശിക്ഷിക്കുക}}
\textbf{\ml{ശിക്ഷിക്കുക}} — \emph{śikṣikkuka} — \textquotedblleft{}to punish\textquotedblright{}.

\begin{quote}
\ml{അധ്യാപകൻ} \ml{ആരെയും} \ml{ശിക്ഷിച്ചില്ല}.
\end{quote}

\begin{quote}
\emph{adhyāpakan āreyuṁ śikṣiccilla.}
\end{quote}

\begin{quote}
The teacher did not punish anyone.
\end{quote}

Past: \ml{ശിക്ഷിച്ചു}, \emph{śikṣiccu}.

The verb of \ml{ശിക്ഷ}: the person punished takes -\ml{നെ}.

\begin{cousinweb}
\textbf{\ml{ശിക്ഷിക്കുക}} (\emph{śikṣikkuka}) is built from Sanskrit \emph{śikṣā} (CDIAL 12432), with the Malayalam verb ending -\ml{ിക്കുക}. It is the verb of \ml{ശിക്ഷ}.
\end{cousinweb}

\begin{grammarlens}[title={What you've built}]
One more word for giving reasons and opinions.
\end{grammarlens}

\subsection*{Your turn}
\begin{itemize}
\raggedright
  \item \emph{Say it:} \emph{śikṣikkuka}
  \item \emph{Say it:} \emph{śikṣikkuka}, once more
  \item \emph{Say it:} \emph{adhyāpakan āreyuṁ śikṣiccilla.} — The teacher did not punish anyone.
  \item \emph{Recall:} say the Malayalam for illegal, then the Malayalam for punishment, then say \emph{śikṣikkuka} again
\end{itemize}

\begin{tcolorbox}[breakable,colback=teal!4,colframe=teal!35!black,title={Before you move on}]
What does \ml{ശിക്ഷിക്കുക} mean? (\textquotedblleft{}To punish\textquotedblright{}.) Say it once more.
\end{tcolorbox}
\section[\texorpdfstring{\ml{വിലക്ക്} (vilakkŭ)}{വിലക്ക് (vilakkŭ)}]{\ml{വിലക്ക്} (vilakkŭ) — a ban}

\label{lesson:ML-C418-vilakku}



\begin{quote}
Before the new one: say the Malayalam for punishment, then the Malayalam for to punish.
\end{quote}

\subsection*{You'll want to know: \ml{വിലക്ക്}}
\textbf{\ml{വിലക്ക്}} — \emph{vilakkŭ} — \textquotedblleft{}a ban\textquotedblright{}.

\begin{quote}
\ml{പുഴയിൽ} \ml{മീൻ} \ml{പിടിക്കാൻ} \ml{ഇപ്പോൾ} \ml{വിലക്ക്} \ml{ഉണ്ട്}.
\end{quote}

\begin{quote}
\emph{puḻayil mīn piṭikkān ippōḷ vilakkŭ uṇṭŭ.}
\end{quote}

\begin{quote}
There is now a ban on fishing in the river.
\end{quote}

The verb is \ml{വിലക്കുക}, to forbid.

\begin{cousinweb}
\textbf{\ml{വിലക്ക്}} (\emph{vilakkŭ}) is native Dravidian: DEDR 5423 lists Malayalam \emph{vilakku}, \textquotedblleft{}a prohibition, an interdict, an embargo\textquotedblright{}, and \emph{vilakkuka}, \textquotedblleft{}to prevent, to prohibit\textquotedblright{}, in an article on stepping aside.
\end{cousinweb}

\begin{grammarlens}[title={What you've built}]
One more word for giving reasons and opinions.
\end{grammarlens}

\subsection*{Your turn}
\begin{itemize}
\raggedright
  \item \emph{Say it:} \emph{vilakkŭ}
  \item \emph{Say it:} \emph{vilakkŭ}, once more
  \item \emph{Say it:} \emph{puḻayil mīn piṭikkān ippōḷ vilakkŭ uṇṭŭ.} — There is now a ban on fishing in the river.
  \item \emph{Recall:} say the Malayalam for punishment, then the Malayalam for to punish, then say \emph{vilakkŭ} again
\end{itemize}

\begin{tcolorbox}[breakable,colback=teal!4,colframe=teal!35!black,title={Before you move on}]
What does \ml{വിലക്ക്} mean? (\textquotedblleft{}A ban\textquotedblright{}.) Say it once more.
\end{tcolorbox}
\section[\texorpdfstring{\ml{നിരോധിക്കുക} (nirōdhikkuka)}{നിരോധിക്കുക (nirōdhikkuka)}]{\ml{നിരോധിക്കുക} (nirōdhikkuka) — to prohibit; to ban}

\label{lesson:ML-C418-nirodhikkuka}



\begin{quote}
Before the new one: say the Malayalam for to punish, then the Malayalam for a ban.
\end{quote}

\subsection*{You'll want to know: \ml{നിരോധിക്കുക}}
\textbf{\ml{നിരോധിക്കുക}} — \emph{nirōdhikkuka} — \textquotedblleft{}to prohibit; to ban\textquotedblright{}.

\begin{quote}
\ml{സർക്കാർ} \ml{പ്ലാസ്റ്റിക്} \ml{കവറുകൾ} \ml{നിരോധിച്ചു}.
\end{quote}

\begin{quote}
\emph{sarkkār plāsṟṟikŭ kavaṟukaḷ nirōdhiccu.}
\end{quote}

\begin{quote}
The government banned plastic bags.
\end{quote}

Past: \ml{നിരോധിച്ചു}, \emph{nirōdhiccu}.

The noun is \ml{നിരോധനം}, a prohibition.

\begin{cousinweb}
\textbf{\ml{നിരോധിക്കുക}} (\emph{nirōdhikkuka}) is built from Sanskrit \emph{nirodha}, \textquotedblleft{}confinement, restraint, suppression\textquotedblright{} (Monier-Williams), with the Malayalam verb ending -\ml{ിക്കുക}.
\end{cousinweb}

\begin{grammarlens}[title={What you've built}]
One more word for giving reasons and opinions.
\end{grammarlens}

\subsection*{Your turn}
\begin{itemize}
\raggedright
  \item \emph{Say it:} \emph{nirōdhikkuka}
  \item \emph{Say it:} \emph{nirōdhikkuka}, once more
  \item \emph{Say it:} \emph{sarkkār plāsṟṟikŭ kavaṟukaḷ nirōdhiccu.} — The government banned plastic bags.
  \item \emph{Recall:} say the Malayalam for to punish, then the Malayalam for a ban, then say \emph{nirōdhikkuka} again
\end{itemize}

\begin{tcolorbox}[breakable,colback=teal!4,colframe=teal!35!black,title={Before you move on}]
What does \ml{നിരോധിക്കുക} mean? (\textquotedblleft{}To prohibit; to ban\textquotedblright{}.) Say it once more.
\end{tcolorbox}
\section[\texorpdfstring{\ml{അനുവദിക്കുക} (anuvadikkuka)}{അനുവദിക്കുക (anuvadikkuka)}]{\ml{അനുവദിക്കുക} (anuvadikkuka) — to allow; to permit}

\label{lesson:ML-C418-anuvadikkuka}



\begin{quote}
Before the new one: say the Malayalam for a ban, then the Malayalam for to prohibit.
\end{quote}

\subsection*{You'll want to know: \ml{അനുവദിക്കുക}}
\textbf{\ml{അനുവദിക്കുക}} — \emph{anuvadikkuka} — \textquotedblleft{}to allow; to permit\textquotedblright{}.

\begin{quote}
\ml{രാത്രി} \ml{പുറത്ത്} \ml{പോകാൻ} \ml{അച്ഛൻ} \ml{അനുവദിച്ചില്ല}.
\end{quote}

\begin{quote}
\emph{rātri puṟattŭ pōkān acchan anuvadiccilla.}
\end{quote}

\begin{quote}
Father did not allow going out at night.
\end{quote}

Past: \ml{അനുവദിച്ചു}, \emph{anuvadiccu}.

What is allowed takes -\ml{ആൻ}: \ml{പോകാൻ} \ml{അനുവദിക്കുക}. The noun is \ml{അനുവാദം}, permission.

\begin{cousinweb}
\textbf{\ml{അനുവദിക്കുക}} (\emph{anuvadikkuka}) is built from Sanskrit \emph{anu-vad}, \textquotedblleft{}to speak after, to repeat\textquotedblright{} (Monier-Williams), with the Malayalam verb ending -\ml{ിക്കുക}; how it came to mean \textquotedblleft{}to allow\textquotedblright{} is not settled.
\end{cousinweb}

\begin{grammarlens}[title={What you've built}]
One more word for giving reasons and opinions.
\end{grammarlens}

\subsection*{Your turn}
\begin{itemize}
\raggedright
  \item \emph{Say it:} \emph{anuvadikkuka}
  \item \emph{Say it:} \emph{anuvadikkuka}, once more
  \item \emph{Say it:} \emph{rātri puṟattŭ pōkān acchan anuvadiccilla.} — Father did not allow going out at night.
  \item \emph{Recall:} say the Malayalam for a ban, then the Malayalam for to prohibit, then say \emph{anuvadikkuka} again
\end{itemize}

\begin{tcolorbox}[breakable,colback=teal!4,colframe=teal!35!black,title={Before you move on}]
What does \ml{അനുവദിക്കുക} mean? (\textquotedblleft{}To allow; to permit\textquotedblright{}.) Say it once more.
\end{tcolorbox}
